\documentclass[11pt]{amsart}

\usepackage[T1]{fontenc}
\usepackage{lmodern}
\usepackage{microtype}
\usepackage{amsmath,amssymb,amsthm}
\usepackage{booktabs}
\usepackage[hidelinks]{hyperref}

\hypersetup{
  pdftitle={A Certified RM(7,3,2) Mixed Radial Moore Graph}
}

\newtheorem{theorem}{Theorem}
\theoremstyle{remark}
\newtheorem{remark}[theorem]{Remark}

\newcommand{\RM}{\mathcal{RM}}
\newcommand{\dist}{\operatorname{dist}}
\newcommand{\ecc}{\operatorname{ecc}}

\title[A certified $\RM(7,3,2)$ graph]
{A Certified $\RM(7,3,2)$ Mixed Radial Moore Graph}
\date{}

\subjclass[2020]{05C12, 05C35, 05C85}
\keywords{
mixed graph, radial Moore graph, degree--diameter problem, certificate
}

\begin{document}

\begin{abstract}
We give an explicit mixed graph on $104$ vertices that is totally
$(7,3)$-regular and has radius $2$ and diameter $3$. Hence an
$\RM(7,3,2)$ mixed radial Moore graph exists. The graph has a unique
central vertex and status $1$-norm $N_1=3378$. We also correct a
single transposed table cell: the value $N_1=18$, correctly attributed
in the text of \cite{CLC2023} to the order-$108$ graphs with
$(r,z)=(3,7)$, is printed in \cite[Table~3]{CLC2023} and
\cite[Table~4.2]{ACL2026} in the cell $(r,z)=(7,3)$, and
\cite[p.~88]{ACL2026} consequently reports the existence of
$\RM(7,3,2)$ as guaranteed. The complete adjacency data of the graph
are printed in Section~1.
\end{abstract}

\maketitle

\section{The construction}

A simple mixed graph $G=(V,E,A)$ has undirected edges
\[
E\subseteq\binom{V}{2}
\]
and directed arcs
\[
A\subseteq\{(u,v)\in V^2:u\ne v\},
\]
with no directed digons and no edge--arc conflicts. Distances respect
arc orientation, and all eccentricities below are out-eccentricities.
The graph is \emph{totally $(r,z)$-regular} if every vertex has
undirected degree $r$ and directed in-degree and out-degree $z$.

For diameter two, the mixed Moore bound is
\begin{equation}\label{eq:moore}
M(r,z,2)
 =1+(r+z)+r(r-1+z)+z(r+z)
 =(r+z)^2+z+1.
\end{equation}
This is the diameter-two mixed Moore bound of Bos\'ak
\cite{Bosak1979}. A totally $(r,z)$-regular mixed graph of order
$M(r,z,k)$, radius $k$, and diameter $k+1$ is an $\RM(r,z,k)$ graph
\cite{CLC2023,ACL2026}.

For a vertex $v$, let
\[
\sigma(v)=\sum_{u\in V}\dist(v,u).
\]
For $(r,z,k)=(7,3,2)$, the Moore status is
\[
\sigma_M=10+2\cdot93=196.
\]
Following \cite{CLC2023}, define
\[
N_1(G)
 =\bigl\|(\sigma(v))_{v\in V}-196\mathbf 1\bigr\|_1.
\]

\begin{theorem}\label{thm:main}
There is an $\RM(7,3,2)$ graph on $104$ vertices with a unique central
vertex and
\[
N_1(G)=3378.
\]
\end{theorem}

\begin{proof}
The adjacency certificate below gives, for every vertex, its undirected
neighbours $U(v)$ and the heads $A^+(v)$ of its outgoing arcs, and so
determines a mixed graph $G$ on
\[
V=\{0,\ldots,103\}.
\]
Undirected adjacency is symmetric there, and there are no loops,
repeated elements, directed digons, or edge--arc conflicts. Every row
lists seven undirected neighbours and three arc heads, and every vertex
occurs as an arc head exactly three times. Hence
\[
|E|=364=\frac{104\cdot7}{2},
\qquad
|A|=312=104\cdot3,
\]
and, at every vertex,
\[
d_E(v)=7,
\qquad
d_A^+(v)=d_A^-(v)=3.
\]

Breadth-first search from every vertex gives
\[
\{\ecc(v):v\in V\}=\{2^1,3^{103}\}.
\]
Thus $G$ has radius $2$, diameter $3$, and a unique central vertex.
The central distance-layer profile is
\[
[1,10,93].
\]
Since
\[
M(7,3,2)=10^2+3+1=104,
\]
this proves that $G$ is an $\RM(7,3,2)$ graph.

It remains to compute the status norm. For $v\in V$, let
\[
q(v)=|\{u\in V:\dist(v,u)=3\}|.
\]
Every vertex has exactly ten vertices at distance one, and all
remaining vertices are at distance two or three. Hence
\[
\sigma(v)
 =10+2(93-q(v))+3q(v)
 =196+q(v).
\]
The total status is
\[
\sum_{v\in V}\sigma(v)=23762.
\]
Therefore
\[
N_1(G)
 =23762-104\cdot196
 =3378.
\]
Equivalently, exactly $3378$ ordered pairs of vertices are at distance
three.
\end{proof}

\paragraph{Verification.}
The structural conditions, the all-pairs distances, the eccentricities,
the statuses, and $N_1$ are finite computations on the $104$ printed
rows, and we performed them by computer directly from that table in
exact integer arithmetic.

\medskip
\noindent\textbf{Adjacency certificate.}
For each vertex $v$, the table gives the undirected neighbours $U(v)$
and the heads $A^+(v)$ of the outgoing arcs. Vertex $0$ is the central
vertex; vertices $1,\ldots,7$ are its undirected neighbours and
$8,9,10$ the heads of its arcs.

\begin{center}
\fontsize{6.8}{7.15}\selectfont
\setlength{\tabcolsep}{3.2pt}
\renewcommand{\arraystretch}{0.96}
\begin{tabular}{rll@{\quad}rll}
\toprule
$v$ & $U(v)$ & $A^+(v)$ & $v$ & $U(v)$ & $A^+(v)$ \\
\midrule
0 & $1,2,3,4,5,6,7$ & $8,9,10$ & 52 & $5,15,33,41,51,72,91$ & $35,43,90$ \\
1 & $0,11,12,13,14,15,16$ & $17,18,19$ & 53 & $20,43,57,61,74,84,101$ & $30,59,98$ \\
2 & $0,20,21,22,23,24,25$ & $26,27,28$ & 54 & $36,49,55,64,65,66,97$ & $25,31,85$ \\
3 & $0,29,30,31,32,33,34$ & $35,36,37$ & 55 & $27,38,54,72,85,100,102$ & $22,80,96$ \\
4 & $0,38,39,40,41,42,43$ & $44,45,46$ & 56 & $6,11,51,61,66,68,81$ & $7,46,88$ \\
5 & $0,47,48,49,50,51,52$ & $53,54,55$ & 57 & $6,16,27,38,46,53,84$ & $2,71,96$ \\
6 & $0,56,57,58,59,60,61$ & $62,63,64$ & 58 & $6,42,47,59,70,88,94$ & $15,71,86$ \\
7 & $0,65,66,67,68,69,70$ & $71,72,73$ & 59 & $6,15,29,31,58,68,90$ & $44,82,91$ \\
8 & $74,75,76,77,78,79,80$ & $81,82,83$ & 60 & $6,21,25,31,34,37,75$ & $14,49,51$ \\
9 & $84,85,86,87,88,89,90$ & $91,92,93$ & 61 & $6,22,39,47,53,56,75$ & $2,55,83$ \\
10 & $94,95,96,97,98,99,100$ & $101,102,103$ & 62 & $12,23,25,28,36,63,85$ & $11,48,74$ \\
11 & $1,26,32,40,56,86,93$ & $2,15,87$ & 63 & $28,29,32,62,70,98,101$ & $15,39,56$ \\
12 & $1,16,18,62,83,90,100$ & $36,74,103$ & 64 & $13,39,54,72,82,95,101$ & $77,79,88$ \\
13 & $1,16,23,64,69,73,79$ & $5,92,93$ & 65 & $7,46,51,54,72,84,96$ & $8,77,97$ \\
14 & $1,17,28,46,67,80,98$ & $7,57,66$ & 66 & $7,54,56,68,74,75,96$ & $6,29,37$ \\
15 & $1,22,39,46,48,52,59$ & $6,30,61$ & 67 & $7,14,25,26,27,40,77$ & $11,87,100$ \\
16 & $1,12,13,28,33,57,95$ & $23,40,61$ & 68 & $7,21,56,59,66,86,103$ & $47,52,65$ \\
17 & $14,32,71,76,80,95,96$ & $28,34,53$ & 69 & $7,13,20,77,82,93,98$ & $3,12,27$ \\
18 & $12,27,33,38,71,73,99$ & $57,64,70$ & 70 & $7,34,37,42,43,58,63$ & $41,92,93$ \\
19 & $26,27,36,77,81,87,91$ & $11,25,34$ & 71 & $17,18,28,35,44,72,73$ & $0,19,46$ \\
20 & $2,44,49,53,69,74,75$ & $21,24,99$ & 72 & $52,55,64,65,71,88,89$ & $10,70,75$ \\
21 & $2,35,60,68,79,88,102$ & $29,31,94$ & 73 & $13,18,22,45,51,71,83$ & $9,21,32$ \\
22 & $2,15,23,37,61,73,84$ & $13,86,97$ & 74 & $8,20,34,53,66,94,99$ & $32,84,96$ \\
23 & $2,13,22,29,42,62,89$ & $41,43,78$ & 75 & $8,20,60,61,66,92,98$ & $9,59,82$ \\
24 & $2,33,34,41,48,76,86$ & $54,80,101$ & 76 & $8,17,24,25,42,51,90$ & $14,69,85$ \\
25 & $2,37,45,60,62,67,76$ & $4,35,63$ & 77 & $8,19,38,47,67,69,92$ & $0,18,32$ \\
26 & $11,19,30,36,67,78,85$ & $4,17,24$ & 78 & $8,26,40,44,80,92,96$ & $10,43,79$ \\
27 & $18,19,45,46,55,57,67$ & $24,75,90$ & 79 & $8,13,21,29,35,89,93$ & $19,37,81$ \\
28 & $14,16,48,62,63,71,82$ & $50,60,72$ & 80 & $8,14,17,35,78,83,93$ & $34,51,99$ \\
29 & $3,23,59,63,79,90,98$ & $67,84,95$ & 81 & $19,36,56,87,89,97,102$ & $17,38,68$ \\
30 & $3,26,32,44,47,84,90$ & $16,60,100$ & 82 & $28,34,45,48,49,64,69$ & $18,30,56$ \\
31 & $3,45,50,59,60,87,96$ & $13,33,42$ & 83 & $12,41,73,80,94,95,101$ & $58,63,76$ \\
32 & $3,11,17,30,45,63,86$ & $12,23,57$ & 84 & $9,22,30,53,57,65,97$ & $21,80,85$ \\
33 & $3,16,18,24,34,38,52$ & $25,26,79$ & 85 & $9,26,41,49,55,62,102$ & $28,38,68$ \\
34 & $3,24,33,60,70,74,82$ & $1,26,39$ & 86 & $9,11,24,32,40,50,68$ & $29,98,99$ \\
35 & $21,39,42,44,71,79,80$ & $77,78,83$ & 87 & $9,19,31,45,81,89,94$ & $20,47,56$ \\
36 & $19,26,54,62,81,91,103$ & $7,67,102$ & 88 & $9,21,58,72,92,94,103$ & $6,38,76$ \\
37 & $22,25,60,70,93,101,103$ & $49,89,100$ & 89 & $9,23,51,72,79,81,87$ & $42,74,97$ \\
38 & $4,18,33,48,55,57,77$ & $58,67,73$ & 90 & $9,12,29,30,48,59,76$ & $20,50,70$ \\
39 & $4,15,35,43,61,64,103$ & $3,62,94$ & 91 & $19,36,41,42,47,52,99$ & $13,90,95$ \\
40 & $4,11,49,50,67,78,86$ & $1,47,81$ & 92 & $50,75,77,78,88,98,101$ & $0,45,58$ \\
41 & $4,24,52,83,85,91,95$ & $8,16,102$ & 93 & $11,37,43,69,79,80,97$ & $39,55,66$ \\
42 & $4,23,35,58,70,76,91$ & $33,51,98$ & 94 & $10,58,74,83,87,88,103$ & $4,23,91$ \\
43 & $4,39,53,70,93,99,100$ & $48,54,101$ & 95 & $10,16,17,41,46,64,83$ & $48,52,59$ \\
44 & $20,30,35,71,78,97,102$ & $40,65,68$ & 96 & $10,17,31,65,66,78,99$ & $5,22,73$ \\
45 & $25,27,31,32,73,82,87$ & $16,20,89$ & 97 & $10,44,50,54,81,84,93$ & $5,14,95$ \\
46 & $14,15,27,57,65,95,102$ & $62,69,94$ & 98 & $10,14,29,63,69,75,92$ & $12,65,66$ \\
47 & $5,30,58,61,77,91,100$ & $41,42,52$ & 99 & $10,18,43,49,74,91,96$ & $40,72,87$ \\
48 & $5,15,24,28,38,82,90$ & $69,76,84$ & 100 & $10,12,43,47,50,55,103$ & $44,53,60$ \\
49 & $5,20,40,54,82,85,99$ & $3,31,64$ & 101 & $37,53,63,64,83,92,102$ & $33,36,103$ \\
50 & $5,31,40,86,92,97,100$ & $22,61,89$ & 102 & $21,44,46,55,81,85,101$ & $45,50,88$ \\
51 & $5,52,56,65,73,76,89$ & $27,49,75$ & 103 & $36,37,39,68,88,94,100$ & $1,78,86$ \\
\bottomrule
\end{tabular}
\end{center}

\section{A parameter correction}

J{\o}rgensen constructed two nonisomorphic mixed Moore graphs of order
$108$ with undirected degree $3$ and directed degree $7$
\cite{Jorgensen2015}. Ceresuela, L\'opez,
and Chemisana applied an arc swap to these graphs and obtained the
status vector
\[
204^{98},205^8,209^2,
\]
which has status $1$-norm
\[
8(205-204)+2(209-204)=18.
\]
Their text identifies the underlying Moore status vector as
\[
\mathbf{s}_{3,7,2}=204^{108},
\]
but Table~3 places the value $18$ in the cell $(r,z)=(7,3)$
\cite[Table~3]{CLC2023}. Abiad, Ceresuela, and L\'opez
subsequently state that the existence of $\RM(7,3,2)$ is guaranteed
and repeat the same assignment in
\cite[p.~88 and Table~4.2]{ACL2026}. This label is contested in print:
Ceresuela's thesis states that J{\o}rgensen's order-$108$
graphs have $r=7$, $z=3$, while it reproduces the status
vector as $\mathbf{s}_{3,7,2}$ \cite{Ceresuela2024}.

The order determines the correct parameter pair in these cases. Put
\[
s=r+z.
\]
By \eqref{eq:moore},
\[
s^2+1\le M(r,z,2)\le s^2+s+1.
\]
If $M(r,z,2)=108$, then $s=10$, and therefore
\[
z=108-10^2-1=7,
\qquad
r=3.
\]
If $M(r,z,2)=104$, then again $s=10$, but
\[
z=104-10^2-1=3,
\qquad
r=7.
\]
Consequently, the published value $N_1=18$ comes from an
$\RM(3,7,2)$ graph and does not provide a certificate for the
transposed parameter pair $\RM(7,3,2)$. Theorem~\ref{thm:main}
supplies such a certificate.

\begin{remark}
The correction concerns the provenance of the entry $18$. It does not
show that $N_1=18$ is impossible for an $\RM(7,3,2)$ graph. No
optimality claim is made for the value $3378$.
\end{remark}

\section*{Data availability}

All data required to verify Theorem~\ref{thm:main} are printed in the
adjacency certificate in Section~1. There are no ancillary files.

\begin{thebibliography}{9}

\bibitem{ACL2026}
A.~Abiad, J.~M. Ceresuela, and N.~L\'opez,
\emph{A note on integer programming methods for mixed radial Moore
graphs},
Discrete Math. Lett. \textbf{17} (2026), 87--92.
\href{https://doi.org/10.47443/dml.2026.107}
{doi:10.47443/dml.2026.107}.

\bibitem{Bosak1979}
J.~Bos\'ak,
\emph{Partially directed Moore graphs},
Math. Slovaca \textbf{29} (1979), no.~2, 181--196.

\bibitem{Ceresuela2024}
J.~M. Ceresuela Torres,
\emph{Graph theory modelling of PV systems \& contributions to the
study of radial Moore graphs},
Ph.D. thesis, Universitat de Lleida, 2024.

\bibitem{CLC2023}
J.~M. Ceresuela, N.~L\'opez, and D.~Chemisana,
\emph{On mixed radial Moore graphs of diameter 3},
Discrete Math. \textbf{346} (2023), 113525.
\href{https://doi.org/10.1016/j.disc.2023.113525}
{doi:10.1016/j.disc.2023.113525}.

\bibitem{Jorgensen2015}
L.~K. J{\o}rgensen,
\emph{New mixed Moore graphs and directed strongly regular graphs},
Discrete Math. \textbf{338} (2015), no.~6, 1011--1016.
\href{https://doi.org/10.1016/j.disc.2015.01.013}
{doi:10.1016/j.disc.2015.01.013}.

\end{thebibliography}

\end{document}